% Part: computability
% Chapter: computability-theory
% Section: normal-form

\documentclass[../../../include/open-logic-section]{subfiles}

\begin{document}

\olfileid{cmp}{thy}{nfm}
\olsection{مبرهنة الصورة السوية}

متى أمكننا وصف تعريف الدالة القابلة للحساب، والتحقق من صحة ما يُدّعى أنه
سجل تام لحسابها عند دخل معين، أمكن أن نسلك سبيلًا لا يختص
بنموذج حساب بعينه. فلطلب قيمة
الدالة القابلة للحساب~${\OLId{latin-ordinary}{f}}$ عند الدخل~${\OLId{latin-ordinary}{x}}$ نجري ما يأتي:
\begin{enumerate}
  \item نفحص الأوصاف الممكنة لسجلات الحساب، واحدًا بعد آخر.
  \item نختبر السجل المعطى: أهو سجل حساب~${\OLId{latin-ordinary}{f}}({\OLId{latin-ordinary}{x}})$؟
  \item فإن وجدنا السجل الصحيح استخرجنا منه قيمة~${\OLId{latin-ordinary}{f}}({\OLId{latin-ordinary}{x}})$.
\end{enumerate}
وصحة هذا المسلك هي مضمون مبرهنة كليني للصورة السوية.

\begin{thm}[مبرهنة كليني للصورة السوية]
\ollabel{thm:normal-form}
توجد علاقة عودية بدائية~${\OLId{latin-looped}{T}}({\OLId{latin-ordinary}{e}},{\OLId{latin-ordinary}{x}},{\OLId{latin-ordinary}{s}})$ ودالة عودية بدائية~${\OLId{latin-looped}{U}}({\OLId{latin-ordinary}{s}})$ لهما
الخاصية الآتية: إذا كانت ${\OLId{latin-ordinary}{f}}$ أي دالة جزئية قابلة للحساب، وجد~${\OLId{latin-ordinary}{e}}$ بحيث
\[
{\OLId{latin-ordinary}{f}}({\OLId{latin-ordinary}{x}}) \simeq {\OLId{latin-looped}{U}}(\umin{{\OLId{latin-ordinary}{s}}}{{\OLId{latin-looped}{T}}({\OLId{latin-ordinary}{e}}, {\OLId{latin-ordinary}{x}}, {\OLId{latin-ordinary}{s}})})
\]
لكل~${\OLId{latin-ordinary}{x}}$.
\end{thm}

\begin{proof}[رسم البرهان]
في أي نموذج حساب نستطيع أن نضبط وصف الدالة القابلة
للحساب~${\OLId{latin-ordinary}{f}}$، ثم نجعل عددًا طبيعيًا~${\OLId{latin-ordinary}{e}}$ رمزًا له. ونضبط كذلك
«متتالية الحساب»، وهي سجل الخطوات التي تُحسب بها الدالة ذات
الفهرس~${\OLId{latin-ordinary}{e}}$ عند الدخل~${\OLId{latin-ordinary}{x}}$، ونجعل للمتتالية عددًا يرمزها~${\OLId{latin-ordinary}{s}}$.
ويُختار هذا الترميز بحيث تكون
\begin{enumerate}
  \item العلاقة ${\OLId{latin-looped}{T}}({\OLId{latin-ordinary}{e}},{\OLId{latin-ordinary}{x}},{\OLId{latin-ordinary}{s}})$، التي تصدق إذا وفقط إذا كان العدد~${\OLId{latin-ordinary}{s}}$
  يرمز متتالية حساب الدالة ذات الفهرس~${\OLId{latin-ordinary}{e}}$ عند الدخل~${\OLId{latin-ordinary}{x}}$، و
  \item الدالة ${\OLId{latin-looped}{U}}({\OLId{latin-ordinary}{s}})$، التي تستخرج من متتالية الحساب ذات الرمز~${\OLId{latin-ordinary}{s}}$
  النتيجة النهائية لذلك الحساب،
\end{enumerate}
قابلتين للحساب. بل إن العلاقة~${\OLId{latin-looped}{T}}$ والدالة~${\OLId{latin-looped}{U}}$ عوديتان بدائيتان.

% تصحيح ترجمة (C231-02): صُرّح بأن U تستخرج النتيجة من سجل الحساب المرمّز.
\textbf{تنبيه تصحيحي على الترجمة الأساس.} أُصلح وصف~${\OLId{latin-looped}{U}}$ ليصرح بأنها
تستخرج النتيجة من متتالية الحساب المرمّزة، لا أنها «تعيّن المتتالية إلى»
نتيجة.
\end{proof}

\begin{explain}
وأما استيفاء البرهان الصارم فيقتضي اختيار نموذج حساب، وبسط
ترميز أوصاف الدوال ومتتاليات حسابها تفصيلًا، ثم إثبات
العودية البدائية للعلاقة~${\OLId{latin-looped}{T}}$ والدالة~${\OLId{latin-looped}{U}}$. غير أن أكثر التطبيقات يكفيها
أن تكون ${\OLId{latin-looped}{T}}$ و~${\OLId{latin-looped}{U}}$ قابلتين للحساب، مع كلية~${\OLId{latin-looped}{U}}$.

ويحسن جمع فكرة البرهان في العبارة الموجزة الآتية:
$\umin{{\OLId{latin-ordinary}{s}}}{{\OLId{latin-looped}{T}}({\OLId{latin-ordinary}{e}},{\OLId{latin-ordinary}{x}},{\OLId{latin-ordinary}{s}})}$ يبحث عن متتالية حساب الدالة ذات الفهرس~${\OLId{latin-ordinary}{e}}$ عند
الدخل~${\OLId{latin-ordinary}{x}}$، و${\OLId{latin-looped}{U}}$ تعيد خرج متتالية الحساب إن أمكن العثور عليها.
\end{explain}

فإذا اخترنا الدوال العودية الجزئية نموذجًا للحساب، كما فعل كليني
أولًا، دلت المبرهنة على كفاية بحث غير محدود واحد،
أي مؤثر واحد $\umin{{\OLId{latin-ordinary}{y}}}{{\OLId{latin-ordinary}{f}}(\vec {\OLId{latin-ordinary}{x}},{\OLId{latin-ordinary}{y}})}$، لتعريف الدالة. ومن هذه
الجهة يقال إن لكل دالة عودية جزئية صورة سوية.

ونستعمل ${\OLId{latin-looped}{T}}$ و${\OLId{latin-looped}{U}}$ في تعريف التعداد $\cfind{{\OLMathNumeral{0}}}$ و$\cfind{{\OLMathNumeral{1}}}$ و
$\cfind{{\OLMathNumeral{2}}}$ و…. ونفترض من الآن اختيارًا ملائمًا
لـ${\OLId{latin-looped}{T}}$ و~${\OLId{latin-looped}{U}}$، ونجعل المعادلة
% تصحيح صياغة كلاسيكية (C231-01): صُرّح بأن الاختيار فرض ثابت من هذا الموضع، لا برهان مؤجل.
\[
\cfind{{\OLId{latin-ordinary}{e}}}({\OLId{latin-ordinary}{x}}) \simeq {\OLId{latin-looped}{U}}(\umin{{\OLId{latin-ordinary}{s}}}{{\OLId{latin-looped}{T}}({\OLId{latin-ordinary}{e}},{\OLId{latin-ordinary}{x}},{\OLId{latin-ordinary}{s}})})
\]
\emph{تعريفًا} لـ$\cfind{{\OLId{latin-ordinary}{e}}}$.

ومن الحقائق التي نحتاج إليها أيضًا:

\begin{thm}
لكل دالة جزئية قابلة للحساب عدد لا نهائي من الفهارس.
\end{thm}

ووجه ذلك بيّن للحدس: إذا أخذنا وصفًا لدالة قابلة للحساب، صنعنا
وصفًا آخر يؤدي أزواج الدخل والخرج نفسها. ولا يلزم إلا أن
نقدّم على الحساب عملًا لا يغيّر نتيجته، كفحص
${\OLMathNumeral{0}}={\OLMathNumeral{0}}$، أو العد إلى ${\OLMathNumeral{5}}$، أو نحو ذلك. فيختلف فهرس الوصف الجديد
عن الأول، مع أن كليهما فهرس للدالة نفسها؛ وإنما اختلفت طريقة
حسابها اختلافًا يسيرًا.

\end{document}
